GNU/Linux is een open source besturingssysteem en niet gebonden aan een bedrijf, maar moest wel zijn plekje in het bedrijfsleven veroveren. Het beheer en het delen van bestanden is misschien wel de belangrijkste functie in de automatisering. Het hele Internet bestaat uit pagina's (bestanden) die informatie delen. Om bestanden te kunnen delen moet je ze eerst kunnen lezen, het was voor GNU/Linux dus ook van belang dat het data van anderen kon lezen en soms ook schrijven. Bedrijven zoals Microsoft en Apple hadden en hebben hun eigen bestandssystemen die vaak niet open source zijn. De ontwikkelaars van Linux moesten dan ook vaak deze bestandssystemen zelf uitvogelen om ze te kunnen ondersteunen in Linux. Daarnaast ontwikkelden ontwikkelaars ook eigen bestandssystemen omdat het altijd beter kan. Het gevolg is dat GNU/Linux heel veel verschillende bestandssystemen ondersteund. GNU/Linux kan disken van Apple en Microsoft Windows lezen en schrijven. Ondanks dat het \textit{ext} filesystem open source is kunnen de commerciele partijen vaak de Linux disken niet lezen.

De belangrijkste bestandssystemen die door GNU/Linux ondersteunt worden zijn:
\begin{description}
\item[ext] Het meest gebruikte bestandssysteem voor GNU/Linux. Het is een open source filesystem. Het kent verschillende versies (ext, ext2, ext3 en ext4), ext3 en ext4 ondersteunen journaling.
\item[fat] is een bestandssysteem van Microsoft dat stamt uit 1977 en in verschillende versies (FAT12, FAT16, FAT32, en exFAT) verbeteringen doorvoerde op het oorspronkelijke concept. Het is uiteindelijk vervangen door NTFS.
\item[vfat] is een extensie boven op een FAT filesystem die het mogelijk maakt om lange filenamen te gebruiken.
\item[ntfs] Het bestandssysteem dat in de 1990s werd ontwikkeld en de standaard werd voor Windows NT systemen met het verschijnen van Windows NT 3.1. Het is een proprietary journaling bestandssysteem.
\item[hfs(+)] HFS en HFS+ zijn de bestandssystemen die door Apple Mac OS werden gebruikt. HFS+, een journaling file system, werd op Mac OS X gebruikt tot 2017
\item[apfs] Sinds 2017 het primaire bestandssysteem voor Mac OS X computers. Voor leesrechten en decryptie is de specificatie beschreven en openbaar, de rest is proprietary. De mogelijkheid om APFS te schrijven is momenteel (2026) nog experimenteel.
\end{description}

